\section*{Chapter \ref{ch:g10:tax-incidence} --- Tax Incidence and Deadweight Loss}
% figdata: figdata/grade-10/10-tax-incidence/out/ (vat_de, vat_steps); ledger g10:de-vat-2020, g10:fns2021, g10:clk2009; Brenland numbers invented.

\begin{solution}{exo:g10:tax-incidence:1}
Sellers keep $15 - 2 = 13$ dalens; the wedge is 2 dalens.
\end{solution}

\begin{solution}{exo:g10:tax-incidence:2}
With $p_s = p_b - 5$: $100 - 2p_b = 3(p_b - 5) - 50$, so $5p_b = 165$, $p_b = 33$, $p_s = 28$ and $Q = 100 - 66 = 34$ (without the tax: 40 units at 30).
\end{solution}

\begin{solution}{exo:g10:tax-incidence:3}
June 2020: 0.8\,\%; July: 0.0\,\%. December 2020: $-0.7$\,\%; January 2021: 1.6\,\%. The largest jumps are from December 2020 to January 2021 ($+2.3$ points), when the old rates returned,
and from June to July 2020 for food ($-3.1$ points, from 4.2 to 1.1\,\%); the food line jumps as much the other way in July 2021
(from 1.2 to 4.3\,\%), a year after the cut.
\end{solution}

\begin{solution}{exo:g10:tax-incidence:4}
Buyers pay 3 more out of 5: 60\,\%. \Cref{prop:g10:tax-incidence:shares} with $b = 2$, $d = 3$: $d/(b + d) = 3/5$.
\end{solution}

\begin{solution}{exo:g10:tax-incidence:5}
$5 \times 34 = 170$.
\end{solution}

\begin{solution}{exo:g10:tax-incidence:6}
$\tfrac12 \times 5 \times (40 - 34) = 15$.
\end{solution}

\begin{solution}{exo:g10:tax-incidence:7}
With $p_b = p_s + 10$: $60 - p_s - 10 = p_s$, so $p_s = 25$, $p_b = 35$, $Q = 25$ (without the tax: 30 at 30). Revenue $10 \times 25 = 250$; \omterm{def:g10:what-prices-do:dwl}{deadweight loss} $\tfrac12 \times 10 \times 5 =
25$.
\end{solution}

\begin{solution}{exo:g10:tax-incidence:8}
$48 - p = 2p - 24$ gives $p^* = 24$, $Q^* = 24$. With the tax, $48 - p_b = 2(p_b - 6) - 24$ gives $p_b = 28$, $p_s = 22$ and $Q = 20$. Buyers bear $4/6 = 2/3$ of the tax, against $1/3$ in
the text: here demand is steeper (1 unit per dalen against 4), so buyers cut their purchases less and bear more.
\end{solution}

\begin{solution}{exo:g10:tax-incidence:9}
$p_b = 1.5\,p_s$: $120 - 6p_s = 2p_s - 24$, so $p_s = 18$, $p_b = 27$, $Q = 12$. Revenue $(27 - 18) \times 12 = 108$.
\end{solution}

\begin{solution}{exo:g10:tax-incidence:10}
$p_s = p_b + 10$: $60 - p_b = p_b + 10$, so $p_b = 25$, $p_s = 35$, $Q = 35$. Cost $10 \times 35 = 350$. \omterm{def:g10:what-prices-do:dwl}{Deadweight loss} $\tfrac12 \times 10 \times (35 - 30) = 25$.
\end{solution}

\begin{solution}{exo:g10:tax-incidence:11}
With 20: $60 - p_b = p_b - 20$, $p_b = 40$, $p_s = 20$, $Q = 20$. \omterm{def:g10:what-prices-do:dwl}{Deadweight loss} $\tfrac12 \times 20 \times 10 = 100$, four times 25; revenue $20 \times 20 = 400$, 1.6 times 250.
\end{solution}

\begin{solution}{exo:g10:tax-incidence:12}
The same tax, collected once from sellers (included in the shelf price) and once from buyers (added at the till), giving different quantities or different prices in all. The study cited
found that showing prices with the tax included cut purchases by 8\,\%: buyers reacted to the same tax differently depending on how it was shown, which the proposition rules out.
\end{solution}

\begin{solution}{exo:g10:tax-incidence:13}
$R(t) = t(Q^* - kt)$ with $k = bd/(b + d)$ is a parabola with zeros at $0$ and $Q^*/k$, so its top is halfway, at $t^* = Q^*/(2k) = Q^*(b + d)/(2bd)$, where $Q = Q^* - kt^* = Q^*/2$. In
\cref{exo:g10:tax-incidence:7}, $b = d = 1$ and $Q^* = 30$: $t^* = 30$, $Q = 15$, revenue $450$ and \omterm{def:g10:what-prices-do:dwl}{deadweight loss} $\tfrac12 \times 30 \times 15 = 225$.
\end{solution}

\begin{solution}{exo:g10:tax-incidence:14}
Supply is \omterm{def:g10:elasticities:perfectly-elastic}{perfectly inelastic}, so landowners bear it all: the rent that users pay does not change, and the owners keep 100 dalens less per hectare. No hectare leaves use, so there is no
\omterm{def:g10:what-prices-do:dwl}{deadweight loss}. Buildings can be built or not: their supply slopes upward, part of a tax on them is passed to tenants, and fewer are built, so there is a \omterm{def:g10:what-prices-do:dwl}{deadweight loss}.
\end{solution}

\begin{solution}{exo:g10:tax-incidence:15}
Buyers compare shelf prices; a tax added at the till is easy to overlook, so they buy as if the price were lower, as the grocery experiment found (showing the tax cut purchases by 8\,\%). The
proposition assumes that buyers react to the full price they pay; if they react only to the shelf price, the quantity depends on whether the tax is in it, and so does who bears it.
\end{solution}

\begin{solution}{pb:g10:tax-incidence:1}
Prices in dalens per litre; quantities in millions of litres; money in millions of dalens a year.

\textbf{1.} $930 - 300p = 600p - 420$ gives $900p = 1350$: $p^* = 1.5$, $Q^* = 480$.

\textbf{2.} $930 - 300p_b = 600(p_b - 0.2) - 420$ gives $900p_b = 1470$, $p_b \approx 1.633$, $p_s \approx 1.433$, and $Q = 930 - 490 = 440$.

\textbf{3.} Drivers pay about 0.133 of the 0.20: two-thirds. \Cref{prop:g10:tax-incidence:shares}: $d/(b + d) = 600/900 = 2/3$.

\textbf{4.} $\varepsilon_d = -300 \times 1.5/480 = -0.9375$, $\varepsilon_s = 600 \times 1.5/480 = 1.875$; share $1.875/(1.875 + 0.9375) = 1.875/2.8125 = 2/3$.

\textbf{5.} Nothing (\cref{prop:g10:tax-incidence:same}): the same prices and quantity.

\textbf{6.} $0.20 \times 440 = 88$.

\textbf{7.} $\tfrac12 \times 0.20 \times (480 - 440) = 4$.

\textbf{8.} $p_b = 1.5 + \tfrac23 \times 0.4 \approx 1.767$, $p_s \approx 1.367$, $Q = 480 - 200 \times 0.4 = 400$. Revenue $0.4 \times 400 = 160$; \omterm{def:g10:what-prices-do:dwl}{deadweight loss} $\tfrac12 \times 0.4 \times 80 = 16$.

\textbf{9.} The \omterm{def:g10:what-prices-do:dwl}{deadweight loss} grew 4 times; the revenue $160/88 \approx 1.8$ times.

\textbf{10.} $Q_t = 480 - 200t$, so $R(t) = t(480 - 200t)$.

\textbf{11.} $t^* = 480/400 = 1.2$ dalens per litre.

\textbf{12.} $Q = 480 - 240 = 240$; $p_b = 1.5 + \tfrac23 \times 1.2 = 2.3$; $p_s = 1.1$.

\textbf{13.} Revenue $1.2 \times 240 = 288$; \omterm{def:g10:what-prices-do:dwl}{deadweight loss} $\tfrac12 \times 1.2 \times 240 = 144$.

\textbf{14.} At 0.20: $4/88 \approx 0.045$. At 0.40: $16/160 = 0.10$. At 1.20: $144/288 = 0.5$ dalen of lost surplus per dalen raised.

\textbf{15.} $1380 - 600 \times 1.5 = 480$, equal to supply at 1.5.

\textbf{16.} $k = 600 \times 600/1200 = 300$: $Q = 480 - 120 = 360$, revenue $0.4 \times 360 = 144$, \omterm{def:g10:what-prices-do:dwl}{deadweight loss} $\tfrac12 \times 0.4 \times 120 = 24$.

\textbf{17.} $d/(b + d) = 600/1200 = 1/2$.

\textbf{18.} $t^* = 480/600 = 0.8$, with $Q = 240$ and revenue $0.8 \times 240 = 192$.

\textbf{19.} When drivers can cut their fuel use more for each dalen of tax, the tax base shrinks faster: each rise of the tax loses more litres. Revenue is smaller at every rate, and the units that
are no longer bought, each carrying surplus, are more numerous.

\textbf{20.} With the demand of Part I, the revenue-maximising tax is 1.20 dalens per litre; it raises 288 million dalens a year and causes a \omterm{def:g10:what-prices-do:dwl}{deadweight loss} of 144 million dalens a year.
\end{solution}
